\documentclass[11pt,a4paper]{article}
\usepackage[utf8]{inputenc}
\usepackage[T1]{fontenc}
\usepackage{amsmath, amssymb, amsthm, bm}
\usepackage{geometry}
\usepackage{xcolor}
\usepackage{tcolorbox}
\usepackage{fancyhdr}

% Define colors and boxes BEFORE babel to avoid active character conflicts
\definecolor{myblue}{RGB}{0, 80, 150}
\definecolor{mygreen}{RGB}{34, 139, 34}

\newtcolorbox{terimbox}[1]{colback=blue!5!white,colframe=myblue,title=#1}
\newtcolorbox{theoriebox}[1]{colback=green!5!white,colframe=mygreen,title=\textbf{Özet ve Püf Noktalar: #1}}
\newtcolorbox{sonucbox}{colback=white,colframe=gray!50}

% Load babel last
\usepackage[turkish]{babel}

\geometry{margin=2.5cm}

\pagestyle{fancy}
\fancyhf{}
\fancyhead[L]{\textbf{Fizik 1 - DM1}}
\fancyhead[R]{\textbf{Konu Anlatımı ve Çözümler}}
\fancyfoot[C]{\thepage}

\title{\textcolor{myblue}{\textbf{Fizik 1 - DM1 Ödev Çözümleri ve Kavramsal Rehber}}}
\author{Hamza}
\date{\today}

\begin{document}

\maketitle

\begin{terimbox}{\textbf{Önemli Fransızca Fizik Terimleri (Terminologie Française)}}
\begin{itemize}
    \item \textbf{Vecteur unitaire:} Birim vektör (uzunluğu 1 olan vektör).
    \item \textbf{Produit scalaire / vectoriel / mixte:} Skaler / Vektörel / Karma çarpım.
    \item \textbf{Repère cartésien / cylindrique:} Kartezyen / Silindirik referans sistemi.
    \item \textbf{Dérivation des vecteurs:} Vektörlerin türevi (zincir kuralına dikkat).
    \item \textbf{Vecteur vitesse:} Hız vektörü (konumun türevi).
    \item \textbf{En déduire:} (Önceki bulgulardan) Yola çıkarak türetmek, çıkarmak.
    \item \textbf{Vérifier:} Doğrulamak, sağlama yapmak.
\end{itemize}
\end{terimbox}

\vspace{0.5cm}

\section*{Exercice 1 : Produits de vecteurs}

\begin{theoriebox}{Vektör Çarpımları ve Bazlar}
\begin{itemize}
    \item \textbf{Norm (Büyüklük):} Bir $\vec{v}$ vektörünün normu $||\vec{v}|| = \sqrt{v_x^2 + v_y^2 + v_z^2}$ formülü ile hesaplanır. Birim vektör ise $\vec{e_v} = \frac{\vec{v}}{||\vec{v}||}$ şeklinde bulunur.
    \item \textbf{Açılar:} İki birim vektörün \textit{skaler çarpımı} aradaki açının kosinüsünü verir ($\vec{e_1} \cdot \vec{e_2} = \cos\theta$). \textit{Vektörel çarpımının} normu ise sinüsünü verir ($||\vec{e_1} \times \vec{e_2}|| = \sin\theta$).
    \item \textbf{Baz Testi:} Üç vektörün bir baz oluşturması için lineer bağımsız olmaları gerekir. Bunun en kolay yolu \textit{determinantın} (karma çarpımın) $0$'dan farklı olduğunu göstermektir.
\end{itemize}
\end{theoriebox}

\subsection*{a) Calcul des normes et des vecteurs unitaires}
Soit les vecteurs $\vec{a}(1,2,2)$, $\vec{b}(2, 2\sqrt{2}, 2)$ et $\vec{c}(0, \sqrt{2}, \sqrt{2})$.
Leurs normes sont :
\begin{align*}
||\vec{a}|| &= \sqrt{1^2 + 2^2 + 2^2} = \sqrt{9} = 3 \\
||\vec{b}|| &= \sqrt{2^2 + (2\sqrt{2})^2 + 2^2} = \sqrt{4 + 8 + 4} = 4 \\
||\vec{c}|| &= \sqrt{0^2 + (\sqrt{2})^2 + (\sqrt{2})^2} = \sqrt{4} = 2
\end{align*}
Les vecteurs unitaires $\vec{e_v} = \frac{\vec{v}}{||\vec{v}||}$ correspondants sont :
\begin{sonucbox}
$\vec{e_a} = \left(\frac{1}{3}, \frac{2}{3}, \frac{2}{3}\right), \quad 
\vec{e_b} = \left(\frac{1}{2}, \frac{\sqrt{2}}{2}, \frac{1}{2}\right), \quad 
\vec{e_c} = \left(0, \frac{\sqrt{2}}{2}, \frac{\sqrt{2}}{2}\right)$
\end{sonucbox}

\subsection*{b) Calcul des cosinus des angles}
Le produit scalaire de deux vecteurs unitaires donne le cosinus de l'angle entre eux :
\begin{align*}
\cos\theta_a &= \vec{e_b} \cdot \vec{e_c} = \left(\frac{1}{2}\right)(0) + \left(\frac{\sqrt{2}}{2}\right)\left(\frac{\sqrt{2}}{2}\right) + \left(\frac{1}{2}\right)\left(\frac{\sqrt{2}}{2}\right) = \frac{2+\sqrt{2}}{4} \\
\cos\theta_b &= \vec{e_a} \cdot \vec{e_c} = \left(\frac{1}{3}\right)(0) + \left(\frac{2}{3}\right)\left(\frac{\sqrt{2}}{2}\right) + \left(\frac{2}{3}\right)\left(\frac{\sqrt{2}}{2}\right) = \frac{4\sqrt{2}}{6} = \frac{2\sqrt{2}}{3} \\
\cos\theta_c &= \vec{e_b} \cdot \vec{e_a} = \left(\frac{1}{2}\right)\left(\frac{1}{3}\right) + \left(\frac{\sqrt{2}}{2}\right)\left(\frac{2}{3}\right) + \left(\frac{1}{2}\right)\left(\frac{2}{3}\right) = \frac{1 + 2\sqrt{2} + 2}{6} = \frac{3+2\sqrt{2}}{6}
\end{align*}

\subsection*{c) Calcul des composantes des vecteurs (produits vectoriels)}
\begin{align*}
\vec{u_a} &= \vec{e_b} \times \vec{e_c} = 
\begin{vmatrix} 
\vec{i} & \vec{j} & \vec{k} \\ 
1/2 & \sqrt{2}/2 & 1/2 \\ 
0 & \sqrt{2}/2 & \sqrt{2}/2 
\end{vmatrix} 
= \vec{i}\left(\frac{2}{4} - \frac{\sqrt{2}}{4}\right) - \vec{j}\left(\frac{\sqrt{2}}{4} - 0\right) + \vec{k}\left(\frac{\sqrt{2}}{4} - 0\right) = \left( \frac{2-\sqrt{2}}{4}, -\frac{\sqrt{2}}{4}, \frac{\sqrt{2}}{4} \right) \\
\vec{u_b} &= \vec{e_c} \times \vec{e_a} = 
\begin{vmatrix} 
\vec{i} & \vec{j} & \vec{k} \\ 
0 & \sqrt{2}/2 & \sqrt{2}/2 \\ 
1/3 & 2/3 & 2/3 
\end{vmatrix} 
= \vec{i}\left(\frac{2\sqrt{2}}{6} - \frac{2\sqrt{2}}{6}\right) - \vec{j}\left(0 - \frac{\sqrt{2}}{6}\right) + \vec{k}\left(0 - \frac{\sqrt{2}}{6}\right) = \left( 0, \frac{\sqrt{2}}{6}, -\frac{\sqrt{2}}{6} \right) \\
\vec{u_c} &= \vec{e_a} \times \vec{e_b} = 
\begin{vmatrix} 
\vec{i} & \vec{j} & \vec{k} \\ 
1/3 & 2/3 & 2/3 \\ 
1/2 & \sqrt{2}/2 & 1/2 
\end{vmatrix} 
= \vec{i}\left(\frac{2}{6} - \frac{2\sqrt{2}}{6}\right) - \vec{j}\left(\frac{1}{6} - \frac{2}{6}\right) + \vec{k}\left(\frac{\sqrt{2}}{6} - \frac{2}{6}\right) = \left( \frac{2-2\sqrt{2}}{6}, \frac{1}{6}, \frac{\sqrt{2}-2}{6} \right)
\end{align*}

\subsection*{d) Déduction des sinus et vérification}
Puisque $||\vec{u}|| = ||\vec{e_1} \times \vec{e_2}|| = ||\vec{e_1}|| \cdot ||\vec{e_2}|| \cdot \sin\theta = \sin\theta$ :
\begin{align*}
\sin\theta_a &= ||\vec{u_a}|| = \sqrt{\left(\frac{2-\sqrt{2}}{4}\right)^2 + \left(-\frac{\sqrt{2}}{4}\right)^2 + \left(\frac{\sqrt{2}}{4}\right)^2} = \sqrt{\frac{10-4\sqrt{2}}{16}} = \frac{\sqrt{10-4\sqrt{2}}}{4} \\
\sin\theta_b &= ||\vec{u_b}|| = \sqrt{0 + \left(\frac{\sqrt{2}}{6}\right)^2 + \left(-\frac{\sqrt{2}}{6}\right)^2} = \sqrt{\frac{4}{36}} = \frac{1}{3} \\
\sin\theta_c &= ||\vec{u_c}|| = \sqrt{\left(\frac{2-2\sqrt{2}}{6}\right)^2 + \left(\frac{1}{6}\right)^2 + \left(\frac{\sqrt{2}-2}{6}\right)^2} = \frac{\sqrt{19-12\sqrt{2}}}{6}
\end{align*}

\textbf{Vérification à l'aide de la question b) :}
On utilise la relation trigonométrique $\sin\theta = \sqrt{1 - \cos^2\theta}$ :
\begin{itemize}
    \item Pour $\theta_a$ : $\cos\theta_a = \frac{2+\sqrt{2}}{4} \implies \cos^2\theta_a = \frac{6+4\sqrt{2}}{16}$. Donc $\sin\theta_a = \sqrt{1 - \frac{6+4\sqrt{2}}{16}} = \frac{\sqrt{10-4\sqrt{2}}}{4}$. (Vérifié)
    \item Pour $\theta_b$ : $\cos\theta_b = \frac{2\sqrt{2}}{3} \implies \cos^2\theta_b = \frac{8}{9}$. Donc $\sin\theta_b = \sqrt{1 - \frac{8}{9}} = \frac{1}{3}$. (Vérifié)
    \item Pour $\theta_c$ : $\cos\theta_c = \frac{3+2\sqrt{2}}{6} \implies \cos^2\theta_c = \frac{17+12\sqrt{2}}{36}$. Donc $\sin\theta_c = \sqrt{1 - \frac{17+12\sqrt{2}}{36}} = \frac{\sqrt{19-12\sqrt{2}}}{6}$. (Vérifié)
\end{itemize}

\subsection*{e) Base et base normée ?}
Pour que les trois vecteurs constituent une base de l'espace, ils doivent être non coplanaires. On vérifie cela en calculant leur produit mixte (le déterminant) :
$$ \det(\vec{u_a}, \vec{u_b}, \vec{u_c}) = 
\begin{vmatrix} 
\frac{2-\sqrt{2}}{4} & -\frac{\sqrt{2}}{4} & \frac{\sqrt{2}}{4} \\[2mm]
0 & \frac{\sqrt{2}}{6} & -\frac{\sqrt{2}}{6} \\[2mm]
\frac{2-2\sqrt{2}}{6} & \frac{1}{6} & \frac{\sqrt{2}-2}{6}
\end{vmatrix} 
$$
En développant par rapport à la première colonne (le terme central est nul car $a_{21}=0$) :
$$ \det = \frac{2-\sqrt{2}}{4}\, M_{11} + \frac{2-2\sqrt{2}}{6}\, M_{31} $$
avec les mineurs :
\begin{align*}
M_{11} &= \left(\frac{\sqrt{2}}{6}\right)\left(\frac{\sqrt{2}-2}{6}\right) - \left(-\frac{\sqrt{2}}{6}\right)\left(\frac{1}{6}\right) = \frac{2-2\sqrt{2}}{36} + \frac{\sqrt{2}}{36} = \frac{2-\sqrt{2}}{36} \\
M_{31} &= \left(-\frac{\sqrt{2}}{4}\right)\left(-\frac{\sqrt{2}}{6}\right) - \left(\frac{\sqrt{2}}{4}\right)\left(\frac{\sqrt{2}}{6}\right) = \frac{2}{24} - \frac{2}{24} = 0
\end{align*}
Donc :
$$ \det = \frac{2-\sqrt{2}}{4} \cdot \frac{2-\sqrt{2}}{36} = \frac{(2-\sqrt{2})^2}{144} = \frac{6-4\sqrt{2}}{144} \neq 0 $$
Puisque le déterminant est non nul, les vecteurs $\vec{u_a}, \vec{u_b}, \vec{u_c}$ \textbf{constituent bien une base}. 
Cependant, pour qu'une base soit normée, il faut que la norme de chaque vecteur soit égale à $1$. Or, d'après la question précédente, $||\vec{u_b}|| = \frac{1}{3} \neq 1$. Cette base \textbf{n'est donc pas normée}.

\newpage

\section*{Exercice 2 : Dérivation des vecteurs unitaires}

\begin{theoriebox}{Vektörlerin Türevi ve Dönen Referans Sistemleri}
\begin{itemize}
    \item \textbf{Zincir Kuralı:} Zamana ($t$) göre türev alırken, açının ($\varphi$) da zamanla değiştiğini hesaba katmalıyız: $\frac{d}{dt} = \frac{d\varphi}{dt} \frac{d}{d\varphi} = \dot{\varphi} \frac{d}{d\varphi}$
    \item \textbf{Kendine Diklik:} Her birim vektörün uzunluğu sabit ($1$) olduğundan, kendi türevi daima kendisine \textit{diktir}. 
    \item \textbf{Bağıl Türev:} Dönen silindirik sistemdeki ($R_{cyl}$) bir gözlemci, kendi eksenlerini ($\vec{e_\rho}, \vec{e_\varphi}$) dönmüyor olarak görür. Bu yüzden o referansta bu vektörlerin türevi $0$'dır.
\end{itemize}
\end{theoriebox}

\subsection*{a) Calcul des dérivées en projection dans la base cartésienne $B$}
Les vecteurs unitaires polaires s'écrivent dans la base cartésienne fixe $B (\vec{e_x}, \vec{e_y}, \vec{e_z})$ :
\begin{align*}
\vec{e_\rho} &= \cos\varphi \vec{e_x} + \sin\varphi \vec{e_y} \\
\vec{e_\varphi} &= -\sin\varphi \vec{e_x} + \cos\varphi \vec{e_y}
\end{align*}
Puisque les vecteurs $\vec{e_x}$ et $\vec{e_y}$ sont fixes dans le repère $R$, on dérive uniquement les coefficients par rapport à $\varphi$ :
\begin{sonucbox}
$$ \left[\frac{d\vec{e_\rho}}{d\varphi}\right]_R = \frac{d(\cos\varphi)}{d\varphi} \vec{e_x} + \frac{d(\sin\varphi)}{d\varphi} \vec{e_y} = -\sin\varphi \vec{e_x} + \cos\varphi \vec{e_y} $$
$$ \left[\frac{d\vec{e_\varphi}}{d\varphi}\right]_R = \frac{d(-\sin\varphi)}{d\varphi} \vec{e_x} + \frac{d(\cos\varphi)}{d\varphi} \vec{e_y} = -\cos\varphi \vec{e_x} - \sin\varphi \vec{e_y} $$
\end{sonucbox}

\subsection*{b) Expressions de ces dérivées dans la base cylindrique $B_{cyl}$}
Pour déduire les expressions dans la base cylindrique, on compare les résultats de la question (a) avec les définitions de $\vec{e_\rho}$ et $\vec{e_\varphi}$.
Pour la première dérivée, on reconnaît exactement l'expression de $\vec{e_\varphi}$ :
$$ -\sin\varphi \vec{e_x} + \cos\varphi \vec{e_y} = \vec{e_\varphi} $$
Pour la deuxième dérivée, en factorisant par le signe moins, on obtient :
$$ -\cos\varphi \vec{e_x} - \sin\varphi \vec{e_y} = -(\cos\varphi \vec{e_x} + \sin\varphi \vec{e_y}) = -\vec{e_\rho} $$
Ainsi, en projection dans la base $B_{cyl}$, on obtient :
\begin{sonucbox}
$$ \left[\frac{d\vec{e_\rho}}{d\varphi}\right]_R = \vec{e_\varphi} $$
$$ \left[\frac{d\vec{e_\varphi}}{d\varphi}\right]_R = -\vec{e_\rho} $$
\end{sonucbox}

\subsection*{c) Expressions par rapport au temps $t$}
Si $\varphi$ est fonction du temps, on applique la règle de la chaîne à l'aide des résultats de la question précédente :
\begin{sonucbox}
$$ \left[\frac{d\vec{e_\rho}}{dt}\right]_R = \left[\frac{d\vec{e_\rho}}{d\varphi}\right]_R \frac{d\varphi}{dt} = \frac{d\varphi}{dt} \vec{e_\varphi} $$
$$ \left[\frac{d\vec{e_\varphi}}{dt}\right]_R = \left[\frac{d\vec{e_\varphi}}{d\varphi}\right]_R \frac{d\varphi}{dt} = -\frac{d\varphi}{dt} \vec{e_\rho} $$
\end{sonucbox}

\subsection*{d) Propriété générale de la dérivée d'un vecteur unitaire}
\textbf{Démonstration :} Soit $\vec{u}$ un vecteur unitaire, donc son produit scalaire par lui-même est constant : $\vec{u} \cdot \vec{u} = 1$.
En dérivant cette relation par rapport au temps :
$$ \frac{d}{dt} (\vec{u} \cdot \vec{u}) = \frac{d\vec{u}}{dt} \cdot \vec{u} + \vec{u} \cdot \frac{d\vec{u}}{dt} = 2 \left( \vec{u} \cdot \frac{d\vec{u}}{dt} \right) = 0 $$
Le produit scalaire étant nul, la dérivée $\frac{d\vec{u}}{dt}$ est orthogonale à $\vec{u}$. Par conséquent, la dérivée de tout vecteur unitaire n'a aucune composante sur lui-même.

\subsection*{e) Le vecteur rotation $\vec{\Omega}$}
On cherche un vecteur $\vec{\Omega} = \Omega_\rho \vec{e_\rho} + \Omega_\varphi \vec{e_\varphi} + \Omega_z \vec{e_z}$ tel que :
$\left[\frac{d\vec{e_\rho}}{dt}\right]_R = \vec{\Omega} \times \vec{e_\rho}$ \quad et \quad $\left[\frac{d\vec{e_\varphi}}{dt}\right]_R = \vec{\Omega} \times \vec{e_\varphi}$ \quad et \quad $\left[\frac{d\vec{e_z}}{dt}\right]_R = \vec{\Omega} \times \vec{e_z}$.

1. Pour $\vec{e_\rho}$ :
\begin{align*}
\vec{\Omega} \times \vec{e_\rho} &= (\Omega_\rho \vec{e_\rho} + \Omega_\varphi \vec{e_\varphi} + \Omega_z \vec{e_z}) \times \vec{e_\rho} = -\Omega_\varphi \vec{e_z} + \Omega_z \vec{e_\varphi}
\end{align*}
Puisque $\left[\frac{d\vec{e_\rho}}{dt}\right]_R = \frac{d\varphi}{dt} \vec{e_\varphi}$, on obtient : $\Omega_\varphi = 0$ et $\Omega_z = \frac{d\varphi}{dt}$.

2. Pour $\vec{e_\varphi}$ :
\begin{align*}
\vec{\Omega} \times \vec{e_\varphi} &= \left(\Omega_\rho \vec{e_\rho} + \frac{d\varphi}{dt} \vec{e_z}\right) \times \vec{e_\varphi} = \Omega_\rho \vec{e_z} - \frac{d\varphi}{dt} \vec{e_\rho}
\end{align*}
Puisque $\left[\frac{d\vec{e_\varphi}}{dt}\right]_R = -\frac{d\varphi}{dt} \vec{e_\rho}$, on obtient : $\Omega_\rho = 0$.

3. Pour $\vec{e_z}$ (Vérification) :
Le vecteur $\vec{e_z}$ est constant dans $R$, donc sa dérivée est nulle.
$\vec{\Omega} \times \vec{e_z} = \left(\frac{d\varphi}{dt} \vec{e_z}\right) \times \vec{e_z} = \vec{0}$. La condition est bien respectée.

On en déduit que le vecteur rotation est :
\begin{sonucbox}
$$ \vec{\Omega} = \frac{d\varphi}{dt} \vec{e_z} $$
\end{sonucbox}

\subsection*{f) Dérivées dans le repère $R_{cyl}$}
\textbf{Pourquoi $\left[\frac{d\vec{e_\rho}}{dt}\right]_{R_{cyl}} = \vec{0}$ et $\left[\frac{d\vec{e_\varphi}}{dt}\right]_{R_{cyl}} = \vec{0}$ ?} \\
Par définition, la notation $\left[\dots\right]_{R_{cyl}}$ implique qu'on dérive par rapport au repère cylindrique. Dans ce repère, qui tourne avec le point, les vecteurs $\vec{e_\rho}$ et $\vec{e_\varphi}$ sont immobiles (constants). La dérivée d'un vecteur constant par rapport à son propre repère est toujours nulle.

\subsection*{g) Calcul des dérivées de $\vec{e_x}$ et $\vec{e_y}$ par rapport à $R_{cyl}$}
On a $\vec{e_x} = \cos\varphi \vec{e_\rho} - \sin\varphi \vec{e_\varphi}$ et $\vec{e_y} = \sin\varphi \vec{e_\rho} + \cos\varphi \vec{e_\varphi}$.
Dans le repère $R_{cyl}$, $\vec{e_\rho}$ et $\vec{e_\varphi}$ sont constants, seul $\varphi$ varie par rapport au temps :
\begin{align*}
\left[\frac{d\vec{e_x}}{dt}\right]_{R_{cyl}} &= -\frac{d\varphi}{dt}\sin\varphi \vec{e_\rho} - \frac{d\varphi}{dt}\cos\varphi \vec{e_\varphi} \\
\left[\frac{d\vec{e_y}}{dt}\right]_{R_{cyl}} &= \frac{d\varphi}{dt}\cos\varphi \vec{e_\rho} - \frac{d\varphi}{dt}\sin\varphi \vec{e_\varphi}
\end{align*}
Ces résultats constituent les \textbf{projections sur la base cylindrique $B_{cyl}$} :
\begin{sonucbox}
$$ \left[\frac{d\vec{e_x}}{dt}\right]_{R_{cyl}} = -\frac{d\varphi}{dt}\sin\varphi \vec{e_\rho} - \frac{d\varphi}{dt}\cos\varphi \vec{e_\varphi} \quad (\text{sur } B_{cyl}) $$
$$ \left[\frac{d\vec{e_y}}{dt}\right]_{R_{cyl}} = \frac{d\varphi}{dt}\cos\varphi \vec{e_\rho} - \frac{d\varphi}{dt}\sin\varphi \vec{e_\varphi} \quad (\text{sur } B_{cyl}) $$
\end{sonucbox}

Ensuite, on factorise ces expressions pour retrouver les vecteurs $\vec{e_x}$ et $\vec{e_y}$ :
Pour $\vec{e_x}$ : $-\frac{d\varphi}{dt} (\sin\varphi \vec{e_\rho} + \cos\varphi \vec{e_\varphi}) = -\frac{d\varphi}{dt} \vec{e_y}$ \\
Pour $\vec{e_y}$ : $\frac{d\varphi}{dt} (\cos\varphi \vec{e_\rho} - \sin\varphi \vec{e_\varphi}) = \frac{d\varphi}{dt} \vec{e_x}$

Ainsi, on obtient les \textbf{projections sur la base cartésienne $B$} :
\begin{sonucbox}
$$ \left[\frac{d\vec{e_x}}{dt}\right]_{R_{cyl}} = -\frac{d\varphi}{dt} \vec{e_y} \quad (\text{sur } B) $$
$$ \left[\frac{d\vec{e_y}}{dt}\right]_{R_{cyl}} = \frac{d\varphi}{dt} \vec{e_x} \quad (\text{sur } B) $$
\end{sonucbox}

\newpage

\section*{Exercice 3 : Vecteur vitesse}

\begin{theoriebox}{Hız Vektörü (Vecteur vitesse)}
\begin{itemize}
    \item Hız, konum vektörünün ($\vec{OP}$) referans sistemine göre zamana göre türevidir.
    \item Eğer bileşenleri tek tek türevliyorsak, \textit{çarpımın türevi} kuralını çok dikkatli uygulamalıyız (hem $\rho$ ve $\varphi$ değişkenlerinin, hem de $\vec{e_\rho}$ vektörünün türevleri alınmalıdır).
\end{itemize}
\end{theoriebox}

\subsection*{a) Expression en coordonnées cartésiennes}
Le vecteur position s'écrit $\vec{OP} = x \vec{e_x} + y \vec{e_y} + z \vec{e_z}$.
Dans le repère $R$, la base $B (\vec{e_x}, \vec{e_y}, \vec{e_z})$ est fixe :
\begin{sonucbox}
$$ \left[\frac{d\vec{OP}}{dt}\right]_R = \frac{dx}{dt} \vec{e_x} + \frac{dy}{dt} \vec{e_y} + \frac{dz}{dt} \vec{e_z} $$
\end{sonucbox}

\subsection*{b) Expression dans $B_{cyl}$ à partir de l'expression cartésienne}
Les relations de passage sont :
$x = \rho \cos\varphi \implies \frac{dx}{dt} = \frac{d\rho}{dt} \cos\varphi - \rho \frac{d\varphi}{dt} \sin\varphi$ \\
$y = \rho \sin\varphi \implies \frac{dy}{dt} = \frac{d\rho}{dt} \sin\varphi + \rho \frac{d\varphi}{dt} \cos\varphi$

En injectant ces expressions et sachant que $\vec{e_x} = \cos\varphi \vec{e_\rho} - \sin\varphi \vec{e_\varphi}$ et $\vec{e_y} = \sin\varphi \vec{e_\rho} + \cos\varphi \vec{e_\varphi}$ :
\begin{align*}
\left[\frac{d\vec{OP}}{dt}\right]_R &= \left(\frac{d\rho}{dt} \cos\varphi - \rho \frac{d\varphi}{dt} \sin\varphi\right)(\cos\varphi \vec{e_\rho} - \sin\varphi \vec{e_\varphi}) \\
&\quad + \left(\frac{d\rho}{dt} \sin\varphi + \rho \frac{d\varphi}{dt} \cos\varphi\right)(\sin\varphi \vec{e_\rho} + \cos\varphi \vec{e_\varphi}) + \frac{dz}{dt} \vec{e_z}
\end{align*}
On développe et on regroupe les termes suivant $\vec{e_\rho}$ et $\vec{e_\varphi}$ :
\textbf{Termes selon $\vec{e_\rho}$ :}
$$ \frac{d\rho}{dt} \cos^2\varphi - \rho \frac{d\varphi}{dt} \sin\varphi \cos\varphi + \frac{d\rho}{dt} \sin^2\varphi + \rho \frac{d\varphi}{dt} \cos\varphi \sin\varphi = \frac{d\rho}{dt}(\cos^2\varphi + \sin^2\varphi) = \frac{d\rho}{dt} $$
\textbf{Termes selon $\vec{e_\varphi}$ :}
$$ -\frac{d\rho}{dt} \cos\varphi \sin\varphi + \rho \frac{d\varphi}{dt} \sin^2\varphi + \frac{d\rho}{dt} \sin\varphi \cos\varphi + \rho \frac{d\varphi}{dt} \cos^2\varphi = \rho \frac{d\varphi}{dt}(\sin^2\varphi + \cos^2\varphi) = \rho \frac{d\varphi}{dt} $$
Finalement :
\begin{sonucbox}
$$ \left[\frac{d\vec{OP}}{dt}\right]_R = \frac{d\rho}{dt} \vec{e_\rho} + \rho \frac{d\varphi}{dt} \vec{e_\varphi} + \frac{dz}{dt} \vec{e_z} $$
\end{sonucbox}

\subsection*{c) Résultat direct depuis $B_{cyl}$}
Le vecteur position s'écrit directement $\vec{OP} = \rho \vec{e_\rho} + z \vec{e_z}$.
En dérivant par rapport au temps dans le repère $R$ :
$$ \left[\frac{d\vec{OP}}{dt}\right]_R = \frac{d\rho}{dt} \vec{e_\rho} + \rho \left[\frac{d\vec{e_\rho}}{dt}\right]_R + \frac{dz}{dt} \vec{e_z} $$
D'après la question 2.c, on sait que $\left[\frac{d\vec{e_\rho}}{dt}\right]_R = \frac{d\varphi}{dt} \vec{e_\varphi}$. Donc :
$$ \left[\frac{d\vec{OP}}{dt}\right]_R = \frac{d\rho}{dt} \vec{e_\rho} + \rho \frac{d\varphi}{dt} \vec{e_\varphi} + \frac{dz}{dt} \vec{e_z} $$
Ce qui est beaucoup plus rapide et direct.

\subsection*{d) Calcul de la vitesse dans le repère $R_{cyl}$}
Si l'on calcule la dérivée de $\vec{OP}$ par rapport au repère cylindrique $R_{cyl}$ :
Les vecteurs unitaires $\vec{e_\rho}$ et $\vec{e_z}$ sont fixes par rapport à ce repère. Ainsi :
\begin{sonucbox}
$$ \left[\frac{d\vec{OP}}{dt}\right]_{R_{cyl}} = \frac{d\rho}{dt} \vec{e_\rho} + \frac{dz}{dt} \vec{e_z} $$
\end{sonucbox}
Le terme angulaire a disparu, il s'agit de la vitesse relative pour un observateur lié au repère tournant.

\end{document}
